% s26_us_draft_v2.tex
% §2.6 Stage B: Profitability Analysis — United States
% Second pass: WLM 3.1 + Jargon Calibration + hetpe-academic-writing
% Explanatory mode on — framework assumed established in §2.3
% Date: 2026-04-06

\subsection{Stage B: Profitability Analysis — United States}
\label{sec:stageB_us}

% -----------------------------------------------------------------------
% OPENING — framework assumed; methodology compressed to 3 sentences
% -----------------------------------------------------------------------

The four-channel decomposition applied here extends the Weisskopf (1979)
framework developed in Section~\ref{sec:framework} by separating nominal
capital productivity into its relative price and real components, so that
the exact identity $\Delta\ln r = \phi^{\mu} + \phi^{Py/PK} + \phi^{B_{r}}
+ \phi^{\pi}$ replaces the standard three-channel form. The capacity
utilization series $\hat{\mu}_{t}$ is the structurally identified
estimate from Stage~A rather than a filter residual, which is what
permits the demand and technology channels to be read as analytically
distinct objects rather than composites of the same identification
artefact. The note on figures preceding this sub-section records the
measurement conventions; what follows is their interpretation.

% Note on figures
\begin{quote}\small
\textit{Note on figures:} All figures cover the US Non-Financial
Corporate Sector, 1940--1978. $r_{t} = \text{GOS}_{t}/\text{KNC}_{t}$;
$\hat{\mu}_{t}$ from Stage~A MPF (pinned at $\hat{\mu}_{1948} = 0.80$);
$B_{r,t} = \text{GVA}_{r,t}/(\text{KNR}_{t}\cdot\hat{\mu}_{t})$;
$Py/PK_{t}$ = GDP deflator / capital price deflator.
Regime shading: Pre-Fordist (1940--1946), Fordist (1947--1973),
Post-Fordist onset (1974--1978). Period windows: WWII buildup
(1940--1944), Interim (1944--1947).
\end{quote}

% -----------------------------------------------------------------------
% MOVEMENT A1 — THE COMPENSATION MECHANISM
% -----------------------------------------------------------------------

The sub-period means in Table~\ref{tab:B3_subperiod} reveal a puzzle
that goes to the core of the Fordist political economy: the nominal
profit rate was essentially flat across the two Fordist sub-periods ---
0.187 in the Early Fordist (1945--1966) and 0.189 in the Late Fordist
(1967--1973) --- even as two of its three constituent channels were
contracting. Mean capacity utilization fell from 0.737 to 0.723 across
the same sub-period boundary, registering the secular deterioration in
the activation of productive capacities that Stage~A's overaccumulation
finding predicts: with the transformation elasticity linking capital
accumulation to productive capacity expansion consistently below unity,
each successive wave of investment expanded the productive frontier more
slowly than it expanded the capital stock, widening the structural gap
between what the economy could produce and what it actually realized.
Mean profit share fell from 0.369 to 0.360, recording the distributional
pressure that the political economy of full employment imposed as the
Fordist settlement reached its later phase. The channel that offset both
declines was real capital productivity at normal capacity, which rose
from 0.586 to 0.635 across the same transition. What this means is that
the Fordist accumulation regime reproduced profitability not by sustaining
demand activation or protecting the profit share but by continuously
raising the real productive efficiency with which each unit of capital
stock generated output at normal utilization --- the technology channel
doing the compensatory work that the demand and distributional channels
could no longer perform. Figures~\ref{fig:B1b} and \ref{fig:B1c} make
this offsetting trajectory visible in indexed form: the $\hat{\mu}$
series declines from its wartime peak toward the Fordist plateau and
continues falling through the 1970s, while the $B_{r}$ series climbs
steadily through the entire Fordist window, the two trajectories moving
in opposite directions around the stable profit rate that their product
sustains.

% -----------------------------------------------------------------------
% MOVEMENT A2 — THE RELATIVE PRICE HEADWIND
% -----------------------------------------------------------------------

Reading the compensation mechanism through quantity channels alone,
however, misses a structural feature that the four-channel decomposition
makes visible and the three-channel framework suppresses: the
output-to-capital relative price was compressing monotonically throughout
the Fordist period, imposing a price-denominated headwind against
profitability that the real capital productivity channel had to overcome
before it could begin to compensate for falling demand activation.
Table~\ref{tab:B3_subperiod} records mean $Py/PK$ falling from 1.369 in
the Pre-Fordist/WWII period to 1.176 in the Early Fordist, 1.140 in the
Late Fordist, and 1.025 at the Post-Fordist onset --- a decline of 25
percentage points across the full window. The implication is that capital
goods prices were rising faster than output prices as a structural feature
of Fordist accumulation dynamics, which means that the nominal rate of
return on each additional unit of capital was being squeezed from the
price side even before the oil shock of 1972--74 made that compression
dramatically visible. The oil shock did not initiate this tendency; it
accelerated a pre-existing secular movement to the point where the price
channel displaced the technology channel as the dominant driver of the
profit rate contraction, a finding to which the swing-level decomposition
will return in Paragraph~C4. For now, the point is that the compensation
mechanism documented in Paragraph~A1 operated within a price environment
that was itself working against profitability recovery --- meaning that
$B_{r}$ had to grow not merely to offset declining $\hat{\mu}$ but to do
so while fighting a structural relative price deterioration that the
three-channel nominal framework attributes entirely to the technology
channel without distinguishing these two qualitatively different sources
of capital productivity pressure.

% -----------------------------------------------------------------------
% MOVEMENT B — TURNING-POINT STRUCTURE
% -----------------------------------------------------------------------

Against this structural background, Figure~\ref{fig:B4} identifies seven
turning points on the nominal profit rate between 1951 and 1974 that
organize the Fordist accumulation dynamics into the swing-level episodes
whose channel composition Table~\ref{tab:B2_peaktotrough} decomposes.
Turning points are identified directly on $r_{t}$ rather than on output
or employment because the object of the decomposition is the profit rate,
not the business cycle, and the two do not always share their timing ---
a profit rate contraction can begin while output is still expanding if
the capital stock is growing faster than output, which is precisely the
structural condition the overaccumulation regime generates. The
seven-swing sequence spans two qualitatively distinct structural phases.
The first four swings (1951--1966) constitute the Fordist settlement's
normal cycle: moderate-amplitude contractions followed by recoveries
that cumulatively lifted the profit rate to its 1966 peak. The remaining
three swings (1966--1974) constitute its unravelling: a deep contraction
that the partial recovery of 1970--1972 could not undo, terminated by
an external shock that compressed what little room for manoeuvre the
institutional settlement retained. The largest expansion in the sample
($\Delta\ln r = +0.336$, 1958--1966) and the most consequential
contraction ($\Delta\ln r = -0.269$, 1966--1970) define the structural
arc; everything else in the decomposition is a commentary on why those
two swings had the channel structure they did.

% -----------------------------------------------------------------------
% MOVEMENT C1 — THE EARLY CONTRACTIONS
% -----------------------------------------------------------------------

The 1951--1954 and 1955--1958 contractions record single-channel
dominance, which is the analytical signature of stagnation tendencies
in Vidal's (2014) sense --- downward pressures on profitability that a
functional accumulation regime can absorb through partial institutional
reconfiguration precisely because only one structural channel is under
serious pressure at a time. In the 1951--1954 contraction ($\Delta\ln r
= -0.133$), the distributional channel accounts for 48.5 percent of the
compounded decline and the demand channel for 41.7 percent, with real
capital productivity (5.7 percent) and the relative price channel (4.1
percent) negligible. The distributional dominance places this episode in
the profit squeeze cell of the Basu (2019) typology: Korean War wage
pressure tightened labour markets and compressed the profit share while
military spending wound down, and the two leading channels dragged
simultaneously without any offsetting force. What makes this a stagnation
tendency rather than a structural crisis, in Vidal's taxonomy, is that
a distributional adjustment --- a restoration of wage restraint as the
wartime labour-market conditions eased --- was structurally sufficient
to address the dominant pressure. The 1955--1958 contraction ($\Delta\ln r
= -0.191$) records a different channel structure but the same
single-channel logic: the demand channel accounts for 60.2 percent of
the compounded decline, with distribution contributing 25.3 percent and
both the technology and relative price channels negligible. This is the
crisis-trigger mechanism identified by \citet{OkishioEtAl2022} most
clearly operative in the Fordist sample --- productive capacities were
in place but could not be activated because effective demand contracted,
and the realization failure is registered in the profit rate through the
$\hat{\mu}$ channel. The single-channel dominance again signals
resolvability: demand stimulus --- which the Eisenhower recession
eventually received through federal spending adjustments --- could in
principle address a $\mu$-dominant contraction without restructuring the
institutional settlement. Both episodes are, in this precise sense,
the normal pathology of a functional Fordist accumulation regime.

% -----------------------------------------------------------------------
% MOVEMENT C2 — THE GREAT FORDIST EXPANSION
% -----------------------------------------------------------------------

The 1958--1966 expansion ($\Delta\ln r = +0.336$) represents the
opposite condition: all four channels contributing positively, no
offsetting tendencies, and the demand channel leading at 47.7 percent.
Understanding why the demand channel led --- and why all other channels
aligned behind it --- requires connecting the accounting to the
institutional history. The accumulation dynamics of this period were
shaped by Vietnam-era military expenditure, which sustained aggregate
demand at a scale sufficient to activate productive capacities and drive
nearly half the profit rate recovery \citep{Baker1996}. The demand
activation operated not merely through the $\hat{\mu}$ channel in
accounting terms but through its structural consequence: when productive
capacities are more fully utilized, real capital productivity at normal
capacity rises mechanically, because the same capital stock is generating
output closer to its frontier performance --- which is why $B_{r}$
contributed 22.4 percent of the expansion even without any acceleration
in the underlying rate of technical change. The still-favourable
output-to-capital price environment (the relative price channel
contributing 13.2 percent) amplified the demand effect through the
nominal profit rate. And the distributional channel contributed the
least of the four (16.7 percent), consistent with the
productivity-indexed wage compromise that the Fordist institutional
settlement had routinized: real wages rose with productivity without
improving the profit share, so distribution remained largely inert
during an expansion in which the other three channels were collectively
doing the work. The 1958--1966 period is thus the Fordist accumulation
regime at its most coherent: the institutional settlement holding,
the demand channel leading, and the structural channels responding
proportionately.

% -----------------------------------------------------------------------
% MOVEMENT C3 — STRUCTURAL CRISIS ONSET 1966-1970
% -----------------------------------------------------------------------

The 1966--1970 contraction is where the decomposition's analytical
architecture pays its fullest dividend, because what the channel shares
record is not merely a large profit rate decline but the specific
structural signature that distinguishes dysfunctionality from stagnation.
No single channel dominates: the profit share accounts for 34.5 percent
of the compounded decline ($\Delta\ln r = -0.269$), real capital
productivity for 29.2 percent, and the demand channel for 29.8 percent,
with the relative price channel contributing a minor 6.5 percent.
All four channels drag simultaneously, and no offsetting force is
operative. To understand why this distribution of channel shares
constitutes a qualitative break rather than merely a larger version of
the earlier contractions, it is necessary to hold the contrast with 1951--1954
and 1955--1958 firmly in view. Those episodes registered one-channel
dominance of between 48 and 60 percent precisely because the Fordist
institutional settlement was doing its job: containing distributional
conflict through productivity-indexed wages, sustaining demand through
fiscal management, and holding the organizational conditions of Taylorist
production sufficiently stable to generate real capital productivity
growth. When one channel came under pressure, the others continued to
support the profit rate, and partial reconfiguration could restore the
dominant channel without touching the others. The 1966--1970 contraction
presents a different structure because the institutional containment
failed simultaneously across all three structural channels, leaving no
channel free to offset while the others adjusted. Three mechanisms
converged to produce this outcome. Full-employment conditions sustained
through military expenditure produced labour-market tightening that gave
the working class bargaining leverage it had not possessed since the
pre-Fordist period, directly eroding the profit share and connecting the
distributional channel to the institutional ceiling of full employment
itself \citep{Baker1996, Vidal2019}. The Taylorist work organisation
that had anchored real capital productivity growth was reaching its
productive limits: the systematic separation of conception from execution
that secured managerial control over the labour process simultaneously
eroded the cooperative and cognitive capacities on which further
productivity improvement depended, as the same deskilling that had served
accumulation in the extensive phase became a fetter on accumulation in
the intensive phase \citep{Braverman1998, Vidal2022}. And the demand
channel weakened as the Vietnam stimulus reached its inflationary ceiling
and the fiscal space for continued expansion contracted. The near-equal
weight of these three mechanisms in the decomposition --- each accounting
for between 29 and 35 percent of the profit rate contraction --- is the
empirical expression of what \citet{Vidal2014} means by the dysfunctionality
of the accumulation regime: not a failure of any single institutional
arrangement, but the exhaustion of the settlement as a whole, manifesting
as simultaneous pressure across every channel the profit rate depends on.

% -----------------------------------------------------------------------
% MOVEMENT C4 — PARTIAL RECOVERY AND OIL SHOCK REATTRIBUTION
% -----------------------------------------------------------------------

The 1970--1972 expansion ($\Delta\ln r = +0.066$) that followed is the
smallest positive swing in the sample and provides a precise reading of
what partial crisis looks like in channel terms. The profit share leads
at 52.7 percent --- the distributional adjustment that the 1966--1970
contraction had compressed was partially reversed as labour-market
pressure eased --- and capacity utilization contributes 34.1 percent,
with real capital productivity at 23.5 percent. What limits the recovery
is visible in the one offsetting channel: the relative price of output
to capital goods contributed $-10.4$ percent, meaning that the secular
compression of $Py/PK$ was already working against profitability recovery
even in a period of distributional adjustment and demand activation.
The circuit was partially restored, but the structural conditions for a
renewed Fordist expansion --- a sustained demand environment, an
organizational basis for productivity growth, a favourable price setting
--- were not simultaneously in place. This is the partial crisis in
Vidal's (2014) sense: resolvable in one channel while the others remain
under structural pressure. The 1972--1974 contraction ($\Delta\ln r =
-0.164$) delivers the analytically most consequential result of the
four-channel decomposition. In the three-channel nominal framework,
this contraction appears technology-led: the nominal capital productivity
channel accounts for the dominant share of the profit rate decline,
which has historically been interpreted as evidence of a structural
deterioration in the productive efficiency of the US capital stock in
the early 1970s. The four-channel decomposition disaggregates that
finding and records a different structure: the relative price channel
accounts for 57.8 percent of the compounded decline, while real capital
productivity is mildly positive, contributing an offsetting $-13.0$
percent of the total swing. What this means is that real productive
efficiency --- the capacity of each unit of real capital to generate
real output at normal utilization --- continued to grow during the oil
shock contraction. The nominal profit rate was compressed not because
the capital stock became less productive in real terms but because the
price of capital goods spiked relative to the price of output, reducing
the nominal rate of return on investment even as the real productive
frontier continued to expand. The oil shock operated through the
$Py/PK$ channel, not through $B_{r}$, and its consequence belongs in the
realization crisis cell of the Basu (2019) typology --- specifically, a
price-mediated realization crisis rather than a demand-quantity collapse
or a structural technology failure.

% -----------------------------------------------------------------------
% MOVEMENT D1 — THE CHANNEL PROGRESSION
% -----------------------------------------------------------------------

Taken together, the four contractions document a structural progression
whose pattern is the analytical core of the present section. The
1951--1954 contraction is $\pi$-led (48.5 percent); the 1955--1958
contraction is $\mu$-led (60.2 percent); the 1966--1970 contraction has
no dominant channel ($\pi$ at 34.5, $B_{r}$ at 29.2, $\mu$ at 29.8);
and the 1972--1974 contraction is $Py/PK$-led (57.8 percent). This
progression maps precisely onto Vidal's (2014) taxonomy of accumulation
regime pathology: single-channel contractions are the empirical form
of stagnation tendencies that a functional regime absorbs; three-channel
reinforcing contractions are the empirical form of structural crisis
that no partial reconfiguration can address; and the price-shock
contraction that terminates the sequence is an external disturbance
that accelerated an institutional collapse already underway rather than
initiating it. The capacity to distinguish between these three categories
--- and to assign the 1972--1974 contraction to the third rather than
the second --- is precisely what the four-channel decomposition enables
and what the three-channel framework suppresses by attributing the
price-mediated contraction to the technology channel. Figure~\ref{fig:B3}
makes this progression visible through the normalized channel shares
across all seven swings: the shift from colour concentration in the early
contractions to near-equal distribution across three channels in the
1966--1970 row, and the subsequent inversion in which green ($B_{r}$)
becomes the sole offsetting bar in 1972--1974, captures the structural
transition in a single figure without requiring the reader to compute
the sub-period averages or interpret the swing shares numerically.

% -----------------------------------------------------------------------
% MOVEMENT D2 — THE LIMIT OF THE COMPENSATION MECHANISM
% -----------------------------------------------------------------------

What the sub-period averages and the swing-level decomposition together
establish is not merely a historical narrative of Fordist profitability
dynamics but a precise account of a structural contradiction: the
mechanism through which the Fordist accumulation regime reproduced
profitability contained within it the conditions that brought that
reproduction to a close. The regime sustained the profit rate at
approximately 0.187--0.189 across the Early and Late Fordist sub-periods
through real capital productivity growth --- $B_{r}$ rising from 0.457
in the Pre-Fordist period to 0.635 in the Late Fordist --- against
simultaneous declines in both the activation of productive capacities
(mean $\hat{\mu}$ falling from 0.875 to 0.723) and the output-to-capital
relative price (mean $Py/PK$ falling from 1.369 to 1.140). The mechanism
was constitutive of the regime rather than incidental to it: it reflects
the overaccumulation character of the US non-financial corporate sector
throughout the Fordist period, with the transformation elasticity below
unity converting each wave of capital accumulation into productive
capacity expansion at a rate that sustained real capital productivity
growth precisely because the frontier expanded more slowly than the stock.
The institutional ceiling of this mechanism is what the 1966--1970
three-channel contraction registers. The demand activation that allowed
$B_{r}$ to rise --- because more fully utilized productive capacities
generate higher capital productivity mechanically, not merely through
technical change --- was simultaneously generating the full-employment
conditions that eroded $\pi$ through labour-market tightening and
the Taylorist organizational pressures that exhausted $B_{r}$ through
the degradation of cooperative productive capacity. The mechanism could
not be sustained precisely because the conditions for its operation
were also the conditions for its dissolution. Demand conditions
intervened on the political economy of the ceiling --- the class struggle
over distribution, utilization, and the institutional forms that govern
whether accumulation can proceed.

% -----------------------------------------------------------------------
% TRANSITION TO §2.7
% -----------------------------------------------------------------------

The profitability analysis has established what channels drove the Fordist
profit rate across sub-periods and swings, and has identified the
1966--1970 three-channel contraction as the structural crisis that
terminated the Fordist compensation mechanism. What it cannot establish
is whether capitalists responded to these channels as independent
signals when deciding how much surplus to reinvest. The decomposition
reveals the accounting structure of profitability; Stage~C asks the
behavioural question that the accounting structure motivates --- whether
the demand channel exerted an independent effect on the recapitalization
rate above and beyond its role through the profit rate, or whether the
Okishio restriction $\beta_{\mu} = \beta_{B_{r}}$ holds, compressing
demand and technology into a single behavioural lever.
